\documentclass{article}
\usepackage{icmlko}
\usepackage{graphicx}
\usepackage{booktabs}
\usepackage{amsmath}

\title{자를 바꾸니 벽이 보였다 \\
  {\large 개입-결과 짝 105개를 세었다. 식별 가능한 것은 0개이고, 막는 것은 결측이 아니다}}

\begin{document}
\maketitle

\begin{abstract}
이 실험실은 900노트 동안 「판 $\rho$」 하나로 판정해 왔다. 그것은 \textbf{순위 예측 정확도}인데
목표는 \textbf{개입 $\to$ 결과} 월드모델이다. 세 사이클(897 n-power · 899 날짜 군집 · 900 정식화)이
같은 벽을 다른 각도로 확인했고 채택 문턱 0.00353 은 12도메인 중 \textbf{0개}를 통과시킨다.
이 노트는 문턱을 낮추지도 판을 개선하지도 않고 \textbf{자를 뗐다.}
대신 잰 것은 「저장소에 개입-결과 짝이 몇 개나 있고 그 짝으로 효과를 식별할 수 있는가」다.
\claimx{짝은 105개이고 식별 가능(A등급)은 \textbf{0개}다.}
{어느 한 짝이라도 식별 5조건을 모두 만족하면 이 문장은 무너진다.
지금은 B 34 · C 71 이고 A 는 없다.}
그리고 막는 것이 뜻밖이었다 --- \textbf{결측이 아니라 「처치 시점이 없다」(77/105)와
「배정 기제가 없다」(105/105)} 이고, 뒤엣것은 \textbf{자료를 더 모아도 안 닫히고 설계로만 닫힌다.}
\end{abstract}

\section{왜 자를 뗐나}

판 $\rho$ 는 12도메인 유보 3{,}775행의 순위 상관이다. 채택 문턱 0.00353 아래에서
오늘도 통과 도메인은 \textbf{0/12} 이고 부족 배수는 5.4$\sim$79.5 다. 세 사이클이 각각
다른 이유로 같은 자리에 섰다: 897 은 잡음 단위가 씨앗이 아니라 날짜 군집이라서,
899 는 축 칸 305개가 모형에 안 닿아서, 900 은 그 칸을 닿게 해도 판이 내려가서.

\claimx{그런데 「판이 나쁘다」와 「자가 그 물음을 못 잰다」를 지금까지 아무도 안 갈랐다.}
{개입 효과를 직접 재는 자를 세워 판 $\rho$ 와 결론이 갈리면 두 진단이 분리된다.
이 노트는 그 자를 세울 수 있는지까지만 본다.}
그래서 이번엔 문턱을 낮추지도 판을 고치지도 않는다. \textbf{다른 것을 잰다} ---
저장소에 「운영자가 고른 개입」과 「그 결과」가 짝지어진 레코드가 몇 개나 있고,
그 짝으로 효과를 \textbf{식별}할 수 있는가.

이 물음은 밖에서 오지 않았다. 노트 900 의 지평 팔이 시장팝업
\texttt{is\_free\_entry}(무료입장) $\cdot$ \texttt{reservation\_required}(예약필수)를
\textbf{결측 문제}로만 다뤘다. 그런데 그 두 열은 \textbf{운영자가 고른 개입 그 자체}이고
결과 변수가 같은 레코드에 있다. 900 은 그것을 「축이 모형에 닿나」로만 물었다.

\section{순서 --- 사전등록이 측정보다 13분 앞선다}

\texttt{docs/prereg\_901\_intervention.md} 는 커밋 \texttt{ae0ed6c49}
(2026-08-11T03:21:42+09:00)로 들어갔고, 측정 산출물 스탬프는 03:34:36(재고)과
03:37:02(식별)이다. 산출물에는 시작$\cdot$끝 시각과 코드$\cdot$입력 sha256 을 박았고
\texttt{git HEAD} 스탬프는 \textbf{쓰지 않았다}(긴 러너에서 그것은 원리상 「시작 시점」이다).

사전등록이 미리 못 박은 것: 「개입」의 조작적 정의(T1 순개입 / T2 조건 / T3 유도),
식별 5조건, 그리고 \claimx{못 잴 경우의 판정문 형태 --- 「못 잔다」가 아니라
\textbf{「무엇이 없어서 못 잰다」}여야 한다.}
{판정문에 「무엇이」가 없으면 이 사전등록의 규칙을 어긴 것이다.}

\section{분모 넷을 먼저 갈랐다}

노트 900 이 실제로 걸린 병이 여기 있다 --- 「MKT 15행이 카테고리로 갈린다(59/82)」가
원천 408 분모와 유보 126 분모를 섞은 문장이었다. 그래서 이번엔 딱지를 넷으로 나눴다:
\textbf{D1 원천 레코드 $\cdot$ D2 축행 $\cdot$ D3 판 채점 유보 $\cdot$ D4 도메인 행.}

\begin{table}[h]\centering\small
\begin{tabular}{lrrrr}
\toprule
도메인 & D1 원천 & D2 축행 & D3 유보 & D4 도메인 \\
\midrule
팝업 & 380 & --- & 65 & 89 \\
시장팝업 & 647 & 249 & 126 & 249 \\
\bottomrule
\end{tabular}
\end{table}

\claimx{팝업 하나에서 분모가 넷 다 다르다.}
{네 수 중 둘이 같아지면 이 구분은 과잉이다. 지금은 380/---/65/89 로 전부 다르다.}

\section{판정 --- A등급 0개}

\begin{figure}[h]\centering
\includegraphics[width=\linewidth]{figs/fig1.pdf}
\caption{(a) 개입-결과 짝 105개의 식별 등급. (b) 무엇이 막나 --- W8(배정 기제 미상)이
105/105, W2 가 105/105, W4(처치 시점 없음)가 77/105. (c) 같은 도메인에서도 분모가 다르다.}
\end{figure}

T1(순개입) 열은 \textbf{12/12 도메인}에 있다. 후보 열 105짝을 식별 5조건에 걸었다:
① 처치/대조가 나뉘나 ② 개입 이전 관측이 있나 ③ 역인과가 아닌가
④ 교란이 관측되나 ⑤ 도메인 안 변이가 있나.

\begin{table}[h]\centering\small
\begin{tabular}{lrl}
\toprule
등급 & 짝 수 & 뜻 \\
\midrule
A & \textbf{0} & 다섯 조건 전부 만족 \\
B & 34 & 일부 만족 \\
C & 71 & 대부분 불만족 \\
\bottomrule
\end{tabular}
\end{table}

사전등록의 예측 넷 중 \textbf{둘이 확인되고 둘이 반증됐다}: P1(T1 열이 6/12 이상)
$\to$ 12/12 확인. P2(처치$\cdot$대조 둘 다 관측되는 짝 3개 이상) $\to$ D1 28짝 $\cdot$
D3 22짝 확인(\textbf{두 수는 분모가 달라 안 섞는다}). \textbf{P3(A등급 1개 이상) $\to$ 반증.}
\textbf{P4(개입 이전 관측이 있는 도메인 0) $\to$ 반증} --- 12/12 에서 있다. 단 \textbf{약한 쪽}
(같은 생산자 단위의 더 이른 결과)이고, \textbf{강한 쪽}(같은 레코드의 개입 이전 결과)은
필드 이름을 전량 훑어 \textbf{0} 이었다.

\section{막는 것은 결측이 아니다}

\claimx{전 도메인 공통으로 막는 것은 둘이고, \textbf{둘 다 결측이 아니다.}}
{어느 한 도메인이라도 처치 시점 기록과 배정 기제 기록을 갖고 있으면 이 문장은 좁혀진다.}

\textbf{W4 --- 처치 시점이 없다}(77/105). 개입 값이 \textbf{언제 정해졌는지} 기록이 없다.
유일한 예외가 팝업 \texttt{intervention.attributes.*} 20열로,
\texttt{ingest/doc\_select.is\_pre\_open()} 이 문서 날짜 $\leq$ 개장일을 \textbf{코드로 강제}한다.
⚠ 다만 밀폐는 아니다 --- 기획서$\cdot$계약서는 날짜를 못 읽으면 통과한다
(\texttt{doc\_select.py:71-79}).

\textbf{W8 --- 배정 기제가 없다}(105/105). 운영자가 왜 그 개입을 골랐는지가 없다.
\claimx{W8 은 자료를 더 모아도 안 닫힌다 --- 설계로만 닫힌다.}
{과거 기록에서 배정 규칙을 복원할 수 있으면 무너진다. 105/105 에서 그런 기록은 0건이다.}

\subsection{노트 900 이 「채우면 열린다」던 그 열}

\begin{figure}[h]\centering
\includegraphics[width=\linewidth]{figs/fig2.pdf}
\caption{(a) 시장팝업 \texttt{is\_free\_entry} --- D3 유보 126 중 84 관측인데
\textbf{대조군이 8행}이다. (b) 지평 접합 깔때기(유보 채점행 분모).}
\end{figure}

\claimx{시장팝업 \texttt{is\_free\_entry} 는 \textbf{채워도 못 잰다.}}
{대조군이 열 단위를 넘기면 무너진다. D3 유보 126 중 관측 84, 그중 대조군 \textbf{8행}이고
D1 274 관측도 True 259 / False 15 다.}
노트 900 의 지평 팔이 이 열의 결측 88행을 「채우면 `BLOCK` 을 뗄 자격이 생긴다」고 적었다.
그 문장은 \textbf{설계 해상도의 자에서는} 옳다. 그러나 \textbf{개입 효과의 자에서는} 채움이
대조군을 만들지 못한다 --- 없는 것은 값이 아니라 \textbf{변이}다.

\section{배선 검사가 세 번 발화했다}

\claimx{초판은 팝업에서 A=6 을 냈다 --- 식별을 \textbf{도메인 단위}로 봤기 때문이다.}
{열 단위 판정이 도메인 단위와 같은 답을 내면 이 수정은 불필요했다. 고치니 A=0 이 됐다.}
태거가 손대지 않은 열이 \textbf{태거의 시간 마스크를 빌려 쓰고 있었다.}

\claimx{영화 D3 되짚기가 \textbf{조용한 0} 이었다.}
{키 정규화를 되돌리면 assert 가 터져야 한다 --- 심어서 확인했다:
\texttt{키 406 $\to$ 매칭 0}.}
축 키 \texttt{KOBIS-20215315} 대 원천 \texttt{20215315} 라 교집합이 비었는데
\textbf{예외가 안 나고 0 이 나왔다.} \texttt{assert} 로 닫고 결함을 심어 발화를 확인했다.
고친 뒤 영화의 사전 관측이 0 $\to$ \textbf{510} 으로 바뀌었다.

\claimx{\textbf{판 유보 3{,}775 는 배선에 따라 다르다.}}
{두 배선이 같은 수를 내면 병기는 불필요하다. 챔피언 \texttt{ff753.shell(base())} = 3{,}775,
생 \texttt{harness.load()} = 3{,}786(팝업 $-6$ $\cdot$ 아이돌 $-26$ $\cdot$ 게임 $+43$).}
앞으로 3{,}775 를 인용할 때는 \textbf{배선을 병기해야 한다.}

W-D 카나리아(없는 열이 「0」이 아니라 「열 없음」으로 나오나)는 12/12 통과했다.

\section{한계 --- 모른다고 적는 것들}

\begin{itemize}
\item \textbf{내 사전등록에 결함이 있었다(자가 적발).} 식1 의 「소수 쪽 10행」이
      \textbf{이진 분할을 전제}해 연속형에 정의가 없다 --- 66짝이 그 때문에 「모른다」로 갔다.
      중앙값 분할 보조 수치를 냈으나 \textbf{판정에는 안 썼다}(사후 규칙 변경이 되므로).
\item \textbf{안 뒤진 원천 셋.} 특히 \texttt{pnl\_monthly.json} 의
      \texttt{pnl\_pre\_total}(오픈 전 판관비 합 214행)은 \textbf{개입 강도}인데
      제외 목록 \texttt{derived.*} 에 걸려 T1 에 안 들어갔다.
      \texttt{lab/popaxes.py:21} 이 *"사전인지 확인이 안 돼 뺀다"* 고 적어 두고 그대로다.
\item \textbf{팝업$\cdot$아이돌 D3 되짚기 불가} --- 배선이 행을 넓힌다.
      「없다」가 아니라 \textbf{「못 셌다」}다.
\item 팝업 \texttt{outcome.daily} 89행의 일별 계열을 안 열었다(패널 후보).
\item \textbf{효과 크기를 일절 안 냈다} --- 설계다. 재고 조사가 먼저다.
\end{itemize}

\section{남는 것 --- 다음 물음은 「무엇을 수집해야 하는가」다}

판을 올려도 W4$\cdot$W8 은 안 닫힌다. 그리고 닫는 데 필요한 것은
\textbf{새 도메인이 아니라 새 필드 둘}이다:

\begin{enumerate}
\item \textbf{개입 값의 결정 시각} --- 그 값이 언제 정해졌는지. W4 를 닫는다
\item \textbf{배정 기제 기록} --- 왜 그 개입을 골랐는지. W8 은 이것 없이는 원리상 안 닫힌다
\end{enumerate}

\claimx{\textbf{⓪-가 에 대한 답: 판 $\rho$ 가 나쁜 게 아니라 자료가 개입 효과를
식별할 꼴이 아니다.}}
{판을 올렸을 때 W4$\cdot$W8 이 닫히면 무너진다. 둘은 모형이 아니라 \textbf{수집}의 성질이다.}

\end{document}
